\documentclass[10pt,a4paper]{amsart}
\usepackage[margin=19mm]{geometry}
\usepackage{amsmath,amssymb,mathtools}
\usepackage[scr=rsfs,cal=euler]{mathalfa}
\usepackage{xfrac}
\usepackage[unicode,hidelinks,bookmarks=false]{hyperref}
\newcommand{\ie}{i.e.\ }
\newcommand{\cf}{cf.\ }
\newcommand{\resp}{respectively\ }
\newcommand{\Subsection}{\subsection}
\providecommand{\nobreakdash}{\nobreak\hbox{-}}
\setlength{\emergencystretch}{2em}
\multlinegap0pt
\usepackage{bm}
\usepackage{bbm}
\usepackage{dsfont}
\usepackage{xspace}
\usepackage{cancel}
\usepackage{mathtools}
\usepackage{amsthm}
\makeatletter
\@ifundefined{theorem}{\newtheorem{theorem}{Theorem}}{}
\@ifundefined{lemma}{\newtheorem{lemma}[theorem]{Lemma}}{}
\makeatother
\title[Regularity theory for degenerate fully nonlinear nonlocal eq]{Regularity theory for degenerate fully nonlinear nonlocal equations with a Hamiltonian term}
\author{Yuzhou Fang, Juha Kinnunen, Chao Zhang}
\date{Source: arXiv:2511.13403v1 (CC BY 4.0)}
\begin{document}
\maketitle

\noindent\emph{Source and licence.} Regularity theory for degenerate fully nonlinear nonlocal equations with a Hamiltonian term. Version arXiv:2511.13403v1 (CC BY 4.0), distributed under the Creative Commons Attribution 4.0 International licence, as recorded in the arXiv licence metadata for this exact version and shown on the abstract page https://arxiv.org/abs/2511.13403v1. Full rights notice: licenses/arxiv-2511.13403-CC-BY-4.0.txt.\par\smallskip

\noindent\textit{Editorial scope.} Counted unit: the Lemma whose source label is \texttt{lem3-2} in the article (source0.tex lines 655--669, source label \texttt{lem3-2}), delivered together with the article's own complete original proof of it (source0.tex lines 671--846, one \texttt{proof} environment of 176 lines). No mathematics is written, repaired or re-derived: statement and proof are the article's own text line for line. Required context retained uncounted: 2-1 (lines 198--201); main2 (lines 246--249); lem2-1 (lines 292--319); lem3-1 (lines 412--419); 3-1-1 (lines 425--430); 3-1-6 (lines 593--596). Every definition, display and label of those lines is retained unchanged. Omitted from this package and not redistributed: 1--197, 202--245, 250--291, 320--411, 420--424, 431--592, 597--654, 670--670, 847--1949. Nothing inside a retained excerpt is omitted or altered: every retained body is delivered exactly as the article's lines are written, so no omission receipt of any kind accompanies this package. Citations: the retained excerpts carry no citation command at all, so no bibliography entry of the article is transcribed and no reference is redistributed. Reference adaptation: none was needed and none was made; every reference inside the delivered unit resolves inside it, so no unresolved reference or citation is produced. Printed slips kept literal: no printed word, symbol, hypothesis or number of the counted statement or of its proof was altered. Layout and class: the article's own class is \texttt{amsart} with packages including CJK, hyperref, cite, amsmath, mathrsfs; the wrapper is the lane's pinned \texttt{amsart} editorial wrapper at 10pt on A4, which restates the theorem-like environments the retained text needs and adds the article's own macro definitions that the retained text needs (none), with extra wrapper packages (bm, bbm, dsfont, xspace, cancel, mathtools). The delivered numbering is the unit's own. Assets: no retained span contains a figure environment, a graphics-inclusion command, a PDF-page inclusion, a table or a TikZ picture; no image file is copied into this package, the package declares no asset dependency and no image file is redistributed with it. Licence: version arXiv:2511.13403v1 of this article is distributed under the Creative Commons Attribution 4.0 International licence, as recorded in the arXiv licence metadata for this exact version and shown on the abstract page https://arxiv.org/abs/2511.13403v1. A full rights notice is delivered at licenses/arxiv-2511.13403-CC-BY-4.0.txt. Characters: every retained excerpt is pure ASCII, so the delivered engine, XeLaTeX with its default Latin Modern text font, reports no missing character.
    \begin{equation}
    \label{2-1}
    |H(x,\xi)|\le \mathcal{M}+\mathcal{H}|\xi|^q, \quad q\ge0,
    \end{equation}

\begin{equation}
\label{main2}
-\Phi(x,Du+\xi)\mathcal{I}_\sigma(u,x)+H(x,Du+\xi)=f(x) \quad\text{in }  B_1, %|Du+\xi|^p
\end{equation}

\begin{lemma}
\label{lem2-1}
Consider the cone $\mathcal{D}=\left\{z\in B_{\delta_0|\overline{a}|}:|\langle\overline{a},z\rangle|\ge(1-\eta_0)|\overline{a}||z|\right\}$ with $\overline{a}=\overline{x}-\overline{y}$, $0<|\overline{a}|\le\frac{1}{8}$ and $\delta_0,\eta_0\in\left(0,\frac{1}{2}\right)$. Let
$$
\delta^2u(x,z)=u(x)-u(x+z)+Du(x)\cdot z, \quad  x,z\in\mathbb{R}^N
$$
for a differentiable function $u:\mathbb{R}^N\rightarrow\mathbb{R}$. Set $\phi(x,y)=\varphi(|x-y|)$.
\begin{itemize}
\item[(1)] If
\begin{equation*}
\varphi(t)=\begin{cases}t-\frac{1}{4}t^{1+\beta}, \quad t\in [0,t_0],\\
\varphi(t_0),  \quad t\in (t_0,\infty),
\end{cases}
\end{equation*}
for some $t_0\in(0,1)$ and $\beta\in(0,1)$, then there is an absolute constant $\varepsilon\in\left(0,\frac{1}{2}\right)$ such that, for $\eta_0,\delta_0\le\min\{\varepsilon, \delta_1|\overline{a}|^\beta\}$ with $\delta_1=\frac{\beta(1+\beta)}{384}$, we obtain
\begin{equation*}
\frac{\beta(1+\beta)}{64}|\overline{a}|^{\beta-1}|z|^2\le \delta^2\phi(\cdot,\overline{y})(\overline{x},z)\le\frac{\beta(1+\beta)}{4}|\overline{a}|^{\beta-1}|z|^2
\end{equation*}
for every $z\in \mathcal{D}$.

\item[(2)] If $\varphi(t)=t^\gamma$ with $\gamma\in(0,1)$, then there exists $\varepsilon\in\left(0,\frac{1}{2}\right)$, depending only on $\gamma$, such that for any $\eta_0,\delta_0\le(0, \varepsilon]$
we infer that
\begin{equation*}
\frac{\gamma(1-\gamma)2^\gamma}{32}|\overline{a}|^{\gamma-2}|z|^2\le \delta^2\phi(\cdot,\overline{y})(\overline{x},z)\le\gamma(1-\gamma)|\overline{a}|^{\gamma-2}|z|^2
\end{equation*}
for every $z\in \mathcal{D}$.
\end{itemize}
\end{lemma}

\begin{lemma}
\label{lem3-1}
Let $\sigma\in(1,2)$ and assume that the conditions $(A_1)$--$(A_3)$ hold with $0\le q\le p$. Let $u\in C(\overline{B}_1)$ be a viscosity solution to \eqref{main2} with $\xi\in\mathbb{R}^N$. Then $u$ is locally Lipschitz continuous in $B_1$, that is, there exists a universal constant $C_{\rm lip}\ge1$, independent of $\xi$, such that
$$
|u(x)-u(y)|\le C_{\rm lip}|x-y|
$$
for every $x,y\in B_{\frac12}$. Furthermore, the $C_{\rm lip}$ is uniformly bounded as $\sigma\rightarrow2$.
\end{lemma}

\begin{equation}
\label{3-1-1}
\psi(x)=m\left[\left(|x|^2-\frac{1}{4}\right)_+\right]^2 \quad\text{and}\quad \omega(t)=\begin{cases}t-\frac{1}{4}t^{1+\beta},\qquad0\le t\le1,\\
\omega(1),\quad t>1,
\end{cases}
\end{equation}

\begin{equation}
\label{3-1-6}
-C(\|f\|_{L^\infty(B_1)}+\|u\|_{L^\infty(B_1)}+\|u\|_{L^1_\sigma}+1)\le J_1+2LI_K[\mathcal{D}\setminus B_\delta](\omega,\overline{a}).
\end{equation}

\begin{lemma}
\label{lem3-2}
Let $\sigma\in(1,2)$ and assume that the conditions $(A_1)$--$(A_3)$ hold with $p< q\le p+1$. Let $u\in C(\overline{B}_1)$ be a viscosity solution to \eqref{main2} with $\xi\in\mathbb{R}^N$.
Then there exists a universal constant $b_0>0$ such that if
\begin{align}
\label{3-2-1}
\mathcal{H}(1+|\xi|^{q-p})\le b
\end{align}
for some $b\le b_0$, then $u$ is locally H\"{o}lder continuous in $B_1$.
In other words, there exist universal constants $\alpha\in(0,1)$ and $C\ge1$ (independent of $\xi$) satisfying
$$
|u(x)-u(y)|\le C|x-y|^\alpha
$$
for every $x,y\in B_{\frac12}$. Furthermore, the constant $C$ is uniformly bounded as $\sigma\rightarrow2$.
\end{lemma}

\begin{proof}
Let
$$
\theta=\left(\frac{b_0}{\mathcal{H}}\right)^\frac{1}{q-p}.
$$
In view of \eqref{3-2-1}, we have
\begin{equation}
\label{3-2-2}
\theta\ge1 \quad\text{and}\quad |\xi|\le \theta.
\end{equation}
The proof is similar to that of Lemma \ref{lem3-1}. We substitute $\omega(t)$ in \eqref{3-1-1} in Lemma \ref{lem3-1} with
$\varphi(t)=t^\alpha$,  $t\in[0,\infty)$ for $\alpha\in(0,1)$ to be fixed. By doubling variables we introduce two functions
\begin{equation*}
\begin{cases}
\widetilde{\phi}(x,y)=L\varphi(|x-y|)+\psi(y),\\
\Psi(x,y)=u(x)-u(y)-\widetilde{\phi}(x,y).
\end{cases}
\end{equation*}
Let $\Psi(\overline{x},\overline{y})=\sup_{B_1\times B_1}\Psi(x,y)$ with $(\overline{x},\overline{y})\in \overline{B}_1\times\overline{B}_1$. We show that$\Psi(\overline{x},\overline{y})\le0$
in order to conclude the local H\"{o}lder continuity.
For a contradiction, assume that $\Psi(\overline{x},\overline{y})>0$.
Then
$$
L|\overline{x}-\overline{y}|^\alpha+m\left[\left(|\overline{y}|^2-\frac{1}{4}\right)_+\right]^2\le2\|u\|_{L^\infty(B_1)}.
$$
We may choose
$$
m>2\left(\frac{16}{5}\right)^2\|u\|_{L^\infty(B_1)}
\quad\text{and}\quad
L>16\|u\|_{L^\infty(B_1)}
$$
such that $|\overline{y}|<\frac{3}{4}$, $|\overline{x}-\overline{y}|<\frac{1}{8}$.
It follows that $(\overline{x},\overline{y})\in B_\frac{7}{8}\times B_\frac{7}{8}\subset B_1\times B_1$ and $\overline{x}\neq\overline{y}$, as well as
\begin{equation}
\label{3-2-3}
|\overline{x}-\overline{y}|\le\left(\frac{2\|u\|_{L^\infty(B_1)}}{L}\right)^\frac{1}{\alpha}.
\end{equation}

As in the proof of Lemma \ref{lem3-1}, we obtain the viscosity inequalities
\begin{equation}
\label{3-2-4}
\begin{cases}
-|D\widetilde{h}(\overline{x})+\xi|^p\mathcal{I}_\sigma(\widetilde{v},\overline{x})+H(\overline{x},D\widetilde{h}(\overline{x})+\xi)\le f(\overline{x}),\\
-|D\widetilde{\eta}(\overline{y})+\xi|^p\mathcal{I}_\sigma(\widetilde{w},\overline{y})+H(\overline{y},D\widetilde{\eta}(\overline{y})+\xi)\ge f(\overline{y}),
\end{cases}
\end{equation}
where
\begin{equation*}
\begin{cases}
\widetilde{h}(x)=\widetilde{\phi}(x,\overline{y})-\widetilde{\phi}(\overline{x},\overline{y})+u(\overline{x}),\\
\widetilde{\eta}(y)=-\widetilde{\phi}(\overline{x},y)+\widetilde{\phi}(\overline{x},\overline{y})+u(\overline{y}),
\end{cases}
\end{equation*}
along with
\begin{equation*}
\widetilde{v}(z)=\begin{cases}\widetilde{h}(z),\quad z\in B_\delta(\overline{x}),\\
u(z)\quad\text{ otherwise},
\end{cases}
\end{equation*}
and
\begin{equation*}
\widetilde{w}(z)=\begin{cases}\widetilde{\eta}(z),\quad z\in B_\delta(\overline{y}),\\
u(z)\quad\text{otherwise},
\end{cases}
\end{equation*}
for $\delta\in(0,\frac18)$. A direct calculation shows that
\begin{equation*}
\begin{cases}
D\widetilde{h}(\overline{x})=\alpha L|\overline{x}-\overline{y}|^{\alpha-2}(\overline{x}-\overline{y}),\\
D\widetilde{\eta}(\overline{y})=-\alpha L|\overline{x}-\overline{y}|^{\alpha-2}(\overline{x}-\overline{y})+4m\left(|\overline{y}|^2-\frac{1}{4}\right)_+\overline{y},
\end{cases}
\end{equation*}
and further
\begin{equation*}
\begin{cases}
|D\widetilde{h}(\overline{x})|=\alpha L|\overline{x}-\overline{y}|^{\alpha-1}, \\
\frac{\alpha L|\overline{x}-\overline{y}|^{\alpha-1}}{2}\le|D\widetilde{\eta}(\overline{y})|\le2\alpha L|\overline{x}-\overline{y}|^{\alpha-1}
\end{cases}
\end{equation*}
for $L\ge\frac{15m}{16\alpha}$. Since
$$
4m\left(|\overline{y}|^2-\frac{1}{4}\right)_+|\overline{y}|\le 4m\frac{5}{16}\cdot\frac{3}{4},$$
we may take $L$ large enough that
$$
\frac{15m}{16}\le\alpha L<\alpha L|\overline{x}-\overline{y}|^{\alpha-1}.
$$
Moreover, by taking $L\ge4\alpha^{-1}\theta$, we obtain
$$
|D\widetilde{h}(\overline{x})+\xi|\le\alpha L|\overline{x}-\overline{y}|^{\alpha-1}+\theta\le 2\alpha L|\overline{x}-\overline{y}|^{\alpha-1}
$$
and
$$
|D\widetilde{h}(\overline{x})+\xi|\ge\alpha L|\overline{x}-\overline{y}|^{\alpha-1}-\theta\ge4\theta|\overline{x}-\overline{y}|^{\alpha-1}-\theta\ge3\theta|\overline{x}-\overline{y}|^{\alpha-1},
$$
where \eqref{3-2-2} and $|\overline{x}-\overline{y}|^{\alpha-1}>1$ were utilized. Analogously, we can treat $|D\widetilde{\eta}(\overline{y})+\xi|$. In summary, we have
\begin{equation}
\label{3-2-5}
1<\theta|\overline{x}-\overline{y}|^{\alpha-1}\le|D\widetilde{h}(\overline{x})+\xi|,|D\widetilde{\eta}(\overline{y})+\xi|\le3\alpha L|\overline{x}-\overline{y}|^{\alpha-1}.
\end{equation}
Merging \eqref{3-2-4}, \eqref{3-2-5} with the condition \eqref{2-1} on $H(\cdot)$ with $p<q\le p+1$, we arrive at
\begin{equation}
\label{3-2-6}
\begin{split}
&\mathcal{I}_\sigma(\widetilde{v},\overline{x})-\mathcal{I}_\sigma(\widetilde{w},\overline{y}) \\
&\qquad\ge-2\|f\|_{L^\infty(B_1)}-\left(\frac{\mathcal{M}+\mathcal{H}|D\widetilde{h}(\overline{x})+\xi|^q}{|D\widetilde{h}(\overline{x})+\xi|^p}+
\frac{\mathcal{M}+\mathcal{H}|D\widetilde{\eta}(\overline{y})+\xi|^q}{|D\widetilde{\eta}(\overline{y})+\xi|^p}\right) \\
&\qquad\ge-2(\mathcal{M}+\|f\|_{L^\infty(B_1)})-\mathcal{H}\left(|D\widetilde{h}(\overline{x})+\xi|^{q-p}+|D\widetilde{\eta}(\overline{y})+\xi|^{q-p}\right) \\
&\qquad\ge-2(\mathcal{M}+\|f\|_{L^\infty(B_1)})-\mathcal{H}(3\alpha L)^{q-p}|\overline{x}-\overline{y}|^{(\alpha-1)(q-p)}.
\end{split}
\end{equation}

The estimate on the term $\mathcal{I}_\sigma(\widetilde{v},\overline{x})-\mathcal{I}_\sigma(\widetilde{w},\overline{y})$ in the previous display is similar to Lemma \ref{lem3-1}
and we obtain
\begin{equation}
\label{3-2-7}
\begin{split}
&-C(\|f\|_{L^\infty(B_1)}+\|u\|_{L^\infty(B_1)}+\|u\|_{L^1_\sigma}+\mathcal{M})\\
&\qquad-\mathcal{H}(3\alpha L)^{q-p}|\overline{x}-\overline{y}|^{(\alpha-1)(q-p)}\le LI_K[\mathcal{D}](\varphi,\overline{a}),
\end{split}
\end{equation}
corresponding to \eqref{3-1-6} as $\delta\rightarrow0$.
Let $l(z)=D_x\varphi(|\overline{x}-\overline{y}|)\cdot z$. There holds
$$
I_K[\mathcal{D}](\varphi,\overline{a})=-\int_{\mathcal{D}}\delta^2\varphi(\cdot,\overline{y})(\overline{x},z)K(z)\,dz
$$
with $\varphi(x,y)=\varphi(|x-y|)$. By virtue of Lemma \ref{lem2-1} (2) and
$$
C_{N,\sigma}\frac{\lambda}{|y|^{N+\sigma}}\le K(y)\le C_{N,\sigma}\frac{\Lambda}{|y|^{N+\sigma}},
$$
we have
\begin{equation}
\label{3-2-8}
\begin{split}
I_K[\mathcal{D}](\varphi,\overline{a})&\le-\frac{\alpha(1-\alpha)2^\alpha C_{N,\sigma}\lambda}{32}|\overline{a}|^{\alpha-2}\int_{\mathcal{D}}
|z|^{2-N-\sigma}\,dz \\
&\le-C|\overline{a}|^{\alpha-2}\eta_0^{\frac{N-1}{2}}(\delta_0|\overline{a}|)^{2-\sigma}
=-\widetilde{C}|\overline{a}|^{\alpha-\sigma},
\end{split}
\end{equation}
where in the last line we have selected $\eta_0=\delta_0=\varepsilon$ with $\varepsilon\in(0,\frac{1}{2})$ a number coming from Lemma \ref{lem2-1} (2), and the constant $\widetilde{C}>0$ is uniformly bounded as $\sigma\rightarrow2$. It follows from \eqref{3-2-7}, \eqref{3-2-8} that
\[
\begin{split}
&C(\|f\|_{L^\infty(B_1)}+\|u\|_{L^\infty(B_1)}+\|u\|_{L^1_\sigma}+\mathcal{M})\\
&\qquad+\mathcal{H}(3\alpha L)^{q-p}|\overline{x}-\overline{y}|^{(\alpha-1)(q-p)}\ge\widetilde{C}L|\overline{x}-\overline{y}|^{\alpha-\sigma}
\end{split}
\]
and further by $0<\alpha<1<\sigma$
\begin{equation}
\label{3-2-9}
\begin{split}
\widetilde{C}L&\le C(\|f\|_{L^\infty(B_1)}+\|u\|_{L^\infty(B_1)}+\|u\|_{L^1_\sigma}+\mathcal{M})\\
&\qquad+\mathcal{H}(3\alpha L)^{q-p}|\overline{x}-\overline{y}|^{(\alpha-1)(q-p)+\sigma-\alpha}
\end{split}
\end{equation}
Observe that, since $q\le p+1$ and $\alpha<1$, we have
$$
(\alpha-1)(q-p)+\sigma-\alpha=(\alpha-1)(q-p-1)+\sigma-1>0.
$$
Therefore, recalling \eqref{3-2-1} and $q-p\le1$, we obtain
$$
\widetilde{C}L\le C(\|f\|_{L^\infty(B_1)}+\|u\|_{L^\infty(B_1)}+\|u\|_{L^1_\sigma}+\mathcal{M})+3b_0\alpha ^{q-p}L,
$$
i.e.,
$$
(\widetilde{C}-3b_0\alpha ^{q-p})L\le C(\|f\|_{L^\infty(B_1)}+\|u\|_{L^\infty(B_1)}+\|u\|_{L^1_\sigma}+\mathcal{M}).
$$
Let
$$
\alpha\le\left(\frac{1}{3}\right)^\frac{1}{q-p} \quad\text{and}\quad  0<b_0\le\frac{\widetilde{C}}{2}.
$$
Then
$$
\frac{\widetilde{C}}{2}L\le C(\|f\|_{L^\infty(B_1)}+\|u\|_{L^\infty(B_1)}+\|u\|_{L^1_\sigma}+\mathcal{M}).
$$
with $C,\widetilde{C}>0$ being two universal constants independent of $\mathcal{M},\mathcal{H}$. Taking sufficiently large $L$ leads to a contradiction from the preceding inequality.
\end{proof}

\end{document}
